\documentclass[11pt,a4paper]{article}

\usepackage[utf8]{inputenc}
\usepackage[T1]{fontenc}
\usepackage[italian]{babel}
\usepackage[a4paper,margin=2.5cm]{geometry}
\usepackage{hyperref}
\usepackage{enumitem}
\usepackage{xcolor}
\usepackage{listings}
\usepackage{titlesec}
\usepackage{parskip}

\hypersetup{
    colorlinks=true,
    linkcolor=violet,
    urlcolor=blue,
    citecolor=black
}

\definecolor{codebg}{RGB}{30,30,40}
\definecolor{codefg}{RGB}{220,220,220}
\definecolor{codekey}{RGB}{180,140,220}
\definecolor{codestr}{RGB}{180,220,120}
\definecolor{codecom}{RGB}{120,140,180}

\lstset{
    basicstyle=\ttfamily\footnotesize\color{codefg},
    backgroundcolor=\color{codebg},
    keywordstyle=\color{codekey}\bfseries,
    stringstyle=\color{codestr},
    commentstyle=\color{codecom}\itshape,
    showstringspaces=false,
    breaklines=true,
    frame=single,
    rulecolor=\color{codebg},
    framerule=0pt,
    xleftmargin=8pt,
    xrightmargin=8pt,
    aboveskip=8pt,
    belowskip=8pt,
    tabsize=2
}

\lstdefinelanguage{tsx}{
    keywords={import, from, export, function, const, let, return, if, else, default, async, await, try, catch, finally, throw, new, true, false, null, undefined, type, interface, useState, useEffect, useNavigate, useParams},
    sensitive=true,
    morecomment=[l]//,
    morecomment=[s]{/*}{*/},
    morestring=[b]',
    morestring=[b]",
    morestring=[b]`
}

\titleformat{\section}{\Large\bfseries\color{violet}}{\thesection.}{0.6em}{}
\titleformat{\subsection}{\large\bfseries}{\thesubsection}{0.5em}{}
\titleformat{\subsubsection}{\normalsize\bfseries\itshape}{}{0pt}{}

\title{\bfseries Riassunto delle lezioni\\[4pt]
\large L7 -- L8 -- L8b -- L8c\\[6pt]
\normalsize Progettazione di applicazioni web full stack}
\author{Corso del Prof.\ Devis Bianchini -- Universit\`a degli Studi di Brescia}
\date{}

\begin{document}
\maketitle

\tableofcontents
\bigskip\hrule\bigskip

% ====================================================================
\section{L7 -- Progettazione model-driven dei contenuti informativi (WebML)}
% ====================================================================

\subsection{Obiettivi della lezione}

La lezione introduce la metodologia \textbf{WebML} (\emph{Web Modeling Language}) per la progettazione di applicazioni web \emph{data-intensive}. L'output atteso \`e duplice:
\begin{itemize}[leftmargin=*]
    \item un modello dei dati progettato secondo i principi della metodologia;
    \item la codifica della logica di backend (con NestJS + TypeORM) guidata da tale modello.
\end{itemize}

\subsection{La metodologia WebML}

WebML \`e una metodologia di sviluppo per applicazioni web \emph{data-intensive}, articolata in due assi principali:
\begin{itemize}[leftmargin=*]
    \item \textbf{Modellazione del contenuto informativo} (entit\`a, propriet\`a, relazioni).
    \item \textbf{Modellazione dell'ipertesto} (struttura delle pagine, navigazione, interazione utente).
\end{itemize}

WebML guida sia la logica di backend sia la logica di frontend e rappresenta un passo avanti rispetto alla classica \emph{Guida di Yale} (Lynch \& Horton, 1999), che si limitava a indicare i passi di strategia, progettazione visiva, accessibilit\`a, usabilit\`a e manutenzione.

Il workflow WebML completo comprende:
\begin{enumerate}[leftmargin=*]
    \item Specifica dei requisiti (gruppi utente, casi d'uso, dizionario dei dati).
    \item Progettazione dei dati.
    \item Progettazione dell'ipertesto.
    \item Implementazione.
    \item Promozione, monitoraggio, valutazione, manutenzione.
\end{enumerate}

\subsection{Specifica dei requisiti}

\textbf{Input:} requisiti di business. \textbf{Output:} specifiche semi-formali, sufficientemente precise ma comprensibili dal cliente. Si divide in due sotto-attivit\`a:

\subsubsection{Raccolta dei requisiti}
Attivit\`a poco strutturata, raccoglie:
\begin{itemize}[leftmargin=*]
    \item utenti e gruppi utente (interni/esterni);
    \item requisiti funzionali (casi d'uso);
    \item requisiti sui dati (dizionario);
    \item mockup dell'interfaccia web;
    \item requisiti non funzionali (usabilit\`a, prestazioni, disponibilit\`a, scalabilit\`a, sicurezza, manutenibilit\`a, multi-device).
\end{itemize}

\subsubsection{Analisi dei requisiti}
Revisione e formalizzazione dei requisiti raccolti tramite linguaggi appropriati:
\begin{itemize}[leftmargin=*]
    \item \textbf{UML} per gruppi di utenti e casi d'uso;
    \item soluzioni ad-hoc per dizionario dati e mockup dell'interfaccia.
\end{itemize}

\subsection{Gerarchia dei gruppi utente}

I gruppi utente si organizzano in gerarchia (es. \emph{User} $\to$ \emph{Registrato}/\emph{NonRegistrato} $\to$ \emph{Admin}/\emph{Manager}/\ldots oppure \emph{Interno}/\emph{Fornitore}/\emph{Cliente}). Va sempre identificato \textbf{almeno un gruppo target} a cui dedicare l'applicazione, per focalizzare meglio le funzionalit\`a. Si sfrutta il \emph{meccanismo di ereditariet\`a} per descrivere casi d'uso comuni.

Per ogni gruppo si compila una scheda con: descrizione, dati di profilo, super-gruppo, sotto-gruppo, oggetti in lettura/scrittura, casi d'uso rilevanti.

\subsection{Requisiti sui dati e dizionario}

Tre domande tipiche guidano la raccolta:
\begin{enumerate}[leftmargin=*]
    \item Quali sono gli oggetti informativi da pubblicare/gestire?
    \item Quali propriet\`a li caratterizzano?
    \item In che modo i vari oggetti sono collegati?
\end{enumerate}

Per ciascun elemento del dizionario si specifica: nome ed eventuali sinonimi, descrizione con istanze di esempio, propriet\`a (attributi), relazioni (incluse generalizzazione/specializzazione), componenti (concetto di \emph{part-of}).

% ====================================================================
\section{L8 -- Sviluppo frontend con React}
% ====================================================================

\subsection{Introduzione e contesto}

\subsubsection{Cos'\`e React.js}
\textbf{React.js} \`e una libreria JavaScript/TypeScript per lo sviluppo front-end di applicazioni web. Sviluppata nei laboratori Facebook nel 2011, usata per la prima volta su Instagram nel 2012, rilasciata pubblicamente nel 2013 con licenza MIT. Oggi \`e utilizzata da aziende come AirBnB, Netflix, Paypal, Uber.

\subsubsection{Approccio SPA}
React si ispira al paradigma delle \textbf{SPA (Single Page Applications)}:
\begin{itemize}[leftmargin=*]
    \item un unico \emph{contenitore} che evolve dinamicamente con i dati dal back-end;
    \item nessun caricamento di nuove pagine dal server a ogni interazione;
    \item sviluppo per \emph{componenti} che interagiscono con le API di React, le quali manipolano il DOM.
\end{itemize}

\subsection{Concetto chiave: i componenti}

In React le applicazioni sono fatte di \textbf{componenti}: oggetti dinamici che ricevono input dall'esterno (interazioni utente, time-out, comandi dal back-end) e producono l'elemento grafico da restituire.

Un componente \`e un ``pezzo'' di UI con:
\begin{itemize}[leftmargin=*]
    \item \textbf{logica};
    \item \textbf{aspetto} (HTML o altri componenti annidati).
\end{itemize}

I componenti possono essere \emph{piccoli} (un button) o \emph{grandi} (intere pagine).

\subsection{File principali di un'app React in Nx}

\subsubsection{\texttt{main.tsx}}
Punto di ingresso che monta l'app nel DOM:

\begin{lstlisting}[language=tsx]
import { StrictMode } from 'react';
import { BrowserRouter } from 'react-router-dom';
import * as ReactDOM from 'react-dom/client';
import App from './app/app';

const root = ReactDOM.createRoot(
  document.getElementById('root') as HTMLElement
);

root.render(
  <StrictMode>
    <BrowserRouter>
      <App />
    </BrowserRouter>
  </StrictMode>
);
\end{lstlisting}

Elementi chiave:
\begin{itemize}[leftmargin=*]
    \item \textbf{StrictMode}: contenitore di sviluppo che segnala pratiche sconsigliate ed effetti collaterali; non aggiunge UI visibile.
    \item \textbf{BrowserRouter}: gestisce la navigazione SPA (cambio di URL come \texttt{/}, \texttt{/about}, \texttt{/login} senza ricaricare la pagina).
    \item \textbf{ReactDOM/client}: API moderna (React 18+) per manipolare il DOM.
    \item \texttt{ReactDOM.createRoot}: cerca in \texttt{index.html} l'elemento con \texttt{id="root"}, lo certifica come \texttt{HTMLElement} e crea una root React moderna.
    \item \texttt{App}: componente principale dell'applicazione (root component).
\end{itemize}

\subsubsection{\texttt{index.html}}
Entry point HTML che contiene il \texttt{<div id="root">} dove React monta i componenti, e lo \texttt{<script type="module" src="/src/main.tsx">}.

\subsubsection{\texttt{app.tsx}}
Componente radice. Importa altri componenti (es. \texttt{NxWelcome}) e li compone con il routing:

\begin{lstlisting}[language=tsx]
import NxWelcome from './nx-welcome';
import { Route, Routes, Link } from 'react-router-dom';

export function App() {
  return (
    <div>
      <NxWelcome title="@org/ui" />
      <Routes>
        <Route path="/" element={<div>Root</div>} />
      </Routes>
    </div>
  );
}
\end{lstlisting}

\textbf{Routes} \`e la parte dinamica: legge l'URL corrente e sostituisce la propria zona con l'HTML/componente indicato nella \texttt{element}. Funziona perch\'e \texttt{App} \`e dentro \texttt{<BrowserRouter>}.

\subsection{Componenti funzionali e prop}

Un componente funzionale \`e una funzione che ritorna JSX, esportata per essere usata altrove:

\begin{lstlisting}[language=tsx]
export function NxWelcome({ title }: { title: string }) {
  return <div>Welcome {title}</div>;
}
export default NxWelcome;
\end{lstlisting}

\begin{itemize}[leftmargin=*]
    \item Le \textbf{prop} sono i dati in input al componente; si destrutturano dall'oggetto argomento.
    \item Si tipizzano con TypeScript (\texttt{title: string}).
    \item \texttt{export default} \`e opzionale, ma comodo quando il file contiene un solo componente: permette \texttt{import NxWelcome from './nx-welcome'} senza parentesi graffe.
\end{itemize}

All'interno del JSX si visualizzano dati con la notazione a graffe \texttt{\{\ldots\}}: prop, valori di stato, variabili locali.

\subsection{La sintassi JSX}

JSX (\emph{JavaScript XML}) \`e HTML scritto all'interno di JavaScript. Permette di posizionare tag direttamente nel DOM senza chiamare \texttt{createElement()} o \texttt{appendChild()}. Regole:
\begin{itemize}[leftmargin=*]
    \item Componenti $\to$ notazione \emph{camel-case} (\texttt{<MyComponent>}).
    \item HTML standard $\to$ minuscolo (\texttt{<div>}).
    \item Ogni tag di apertura deve avere uno di chiusura.
\end{itemize}

\subsection{Rendering condizionale}

La UI cambia in base a una condizione/stato:

\begin{lstlisting}[language=tsx]
let content;
if (isLoggedIn) {
  content = <AdminPanel />;
} else {
  content = <LoginForm />;
}
return <div>{content}</div>;
\end{lstlisting}

\subsection{Liste}

Si usa \texttt{map()} per generare elementi; ogni elemento richiede una \textbf{key} univoca:

\begin{lstlisting}[language=tsx]
const listItems = products.map(product =>
  <li
    key={product.id}
    style={{ color: product.isFruit ? 'magenta' : 'darkgreen' }}
  >
    {product.title}
  </li>
);
return <ul>{listItems}</ul>;
\end{lstlisting}

\subsection{Eventi}

React gestisce eventi come click, double-click, input. La sintassi \`e simile a JavaScript ma in JSX (camel-case):

\begin{lstlisting}[language=tsx]
function MyButton() {
  function handleClick() {
    alert('You clicked me!');
  }
  return <button onClick={handleClick}>Click me</button>;
}
\end{lstlisting}

\subsection{Stato (state)}

Lo stato \`e interno al componente e serve a:
\begin{itemize}[leftmargin=*]
    \item aggiornare la UI;
    \item reagire alle interazioni.
\end{itemize}

Quando lo stato cambia, React aggiorna la UI di conseguenza. Si usa l'hook \texttt{useState}:

\begin{lstlisting}[language=tsx]
import { useState } from 'react';

function MyButton() {
  const [count, setCount] = useState(0);
  function handleClick() {
    setCount(count + 1);
  }
  return (
    <button onClick={handleClick}>
      Clicked {count} times
    </button>
  );
}
\end{lstlisting}

\subsection{Condividere i dati tra componenti}

I dati possono essere:
\begin{itemize}[leftmargin=*]
    \item \textbf{passati dall'alto} come prop;
    \item \textbf{condivisi} sollevando lo stato al genitore comune (\emph{lifting state up}).
\end{itemize}

\subsection{Componenti e stile}

Cinque approcci per dare uno stile ai componenti React:

\subsubsection{1. Stile inline}
Veloce ma poco leggibile. L'attributo \texttt{style=\{\{ \ldots \}\}} accetta un oggetto JavaScript con propriet\`a in camel-case (\texttt{font-weight} $\to$ \texttt{fontWeight}), valori come stringhe.

\subsubsection{2. File CSS classico}
File \texttt{.css} separato importato con \texttt{import './style.css'}; nei componenti si usa \texttt{className} al posto di \texttt{class}. Semplice e standard, ma soggetto a conflitti di nome globali.

\subsubsection{3. CSS Modules}
File \texttt{.module.css} associato al componente. Approccio molto usato in Nx, niente conflitti di nome:

\begin{lstlisting}[language=tsx]
import styles from './nx-welcome.module.css';

<div className={styles.welcome}>...</div>
\end{lstlisting}

\subsubsection{4. Fogli di stile globali}
Linkati direttamente in \texttt{index.html} con \texttt{<link rel="stylesheet">}, validi per tutti i componenti, classi referenziate via \texttt{className}.

\subsubsection{5. \texttt{dangerouslySetInnerHTML}}
Approccio ``non standard'' e pericoloso. Forza React ad accettare HTML non filtrato (es. tag \texttt{<style>} inline). Esposto a vulnerabilit\`a \textbf{XSS} (cross-site scripting); da evitare salvo necessit\`a.

\subsection{Dietro le quinte: il Virtual DOM}

Quando si verifica un evento, React \emph{non} manipola direttamente il DOM del browser come fanno le librerie tradizionali, ma manipola un \textbf{Virtual DOM}: una copia esatta del DOM mantenuta in memoria centrale. Pi\`u leggera da modificare, viene poi confrontata (operazione di \emph{diffing}) con la copia precedente. React invia al DOM del browser \emph{solo le modifiche strettamente necessarie}.

\subsubsection{Esempio concreto}
Lista di 10 check-box, utente clicca il primo:
\begin{itemize}[leftmargin=*]
    \item framework tradizionale: ri-renderizza tutti i 10 elementi.
    \item React: modifica solo l'elemento cambiato nel Virtual DOM, fa diffing, e invia al DOM del browser solo l'update del primo check-box (\emph{rendering selettivo}).
\end{itemize}

\subsection{Vantaggi di React}

\textbf{Per lo sviluppatore.} L'approccio a componenti consente riuso e composizione di UI complesse da semplici \emph{mattoncini}; lo sviluppatore si concentra sulla logica e sulla disposizione, mentre Virtual DOM e comunicazione col browser sono interamente a carico di React.

\textbf{Per l'utente.} Il rendering selettivo alleggerisce il processo di aggiornamento dell'interfaccia, con prestazioni percepibili migliori.

\subsection{Riferimenti utili}
\begin{itemize}[leftmargin=*]
    \item Documentazione ufficiale: \url{https://react.dev/learn}
    \item Convertitore HTML$\to$JSX: \url{https://transform.tools/html-to-jsx}
\end{itemize}

% ====================================================================
\section{L8b -- Laboratorio React (catalogo libri + autenticazione)}
% ====================================================================

\subsection{Obiettivo della lezione}

Creazione di una pagina front-end React per visualizzare un catalogo di libri, mettendo in comunicazione frontend e backend; la pagina \`e protetta da meccanismi di autenticazione e autorizzazione. Output: esempi di pagine React con relative rotte e un meccanismo di interazione completo tra le due tier.

\subsection{Procedura}

\begin{enumerate}[leftmargin=*]
    \item Predisporre una funzionalit\`a di backend per estrarre il catalogo.
    \item Abilitare la comunicazione tra frontend e backend (CORS).
    \item Predisporre rotte e pagine nell'applicazione frontend.
\end{enumerate}

\subsection{Back-end: estrazione del catalogo}

\subsubsection{Service}
Il file \texttt{books.service.ts} interroga il repository TypeORM e \textbf{mappa l'entit\`a in un DTO ``flat'' di interfaccia} adatto al frontend:

\begin{lstlisting}[language=tsx]
async findAllBooks(): Promise<BookListItem[]> {
  const books = await this.bookRepository.find({
    relations: { authors: true, category: true },
    order: { title: 'ASC' }
  });

  return books.map((book) => ({
    id: book.id,
    title: book.title,
    publishedYear: book.publishedYear,
    category: { id: book.category.id, name: book.category.name },
    authors: book.authors.map((author) => ({
      id: author.id,
      firstName: author.firstName,
      lastName: author.lastName
    }))
  }));
}
\end{lstlisting}

\subsubsection{Interfaccia condivisa}
\begin{lstlisting}[language=tsx]
export interface BookListItem {
  id: number;
  title: string;
  publishedYear: number;
  category: { id: number; name: string; };
  authors: {
    id: number;
    firstName: string;
    lastName: string;
  }[];
}
\end{lstlisting}

\subsubsection{Controller}
Espone la rotta protetta con \texttt{JwtAuthGuard + RolesGuard} e ritorna il DTO:

\begin{lstlisting}[language=tsx]
@ApiTags('Library APIs')
@Controller('books')
export class OrgBooksController {
  constructor(private orgBooksService: OrgBooksService) {}

  @Get()
  @UseGuards(JwtAuthGuard, RolesGuard)
  @Roles(UserRole.ADMIN)
  @ApiBearerAuth()
  findAll(): Promise<BookListItem[]> {
    return this.orgBooksService.findAllBooks();
  }
}
\end{lstlisting}

\subsection{Cross-Origin Resource Sharing (CORS)}

Per permettere a React (su \texttt{localhost:4200}) di chiamare NestJS (su \texttt{localhost:3333}) servono header CORS espliciti. Nel \texttt{main.ts} del backend:

\begin{lstlisting}[language=tsx]
async function bootstrap() {
  const app = await NestFactory.create(AppModule);

  app.enableCors({
    origin: 'http://localhost:4200',
    credentials: true,
    allowedHeaders: ['Content-Type', 'Authorization'],
    methods: ['GET', 'POST', 'PUT', 'PATCH', 'DELETE', 'OPTIONS']
  });

  const globalPrefix = 'api';
  app.setGlobalPrefix(globalPrefix);
}
\end{lstlisting}

La configurazione permette:
\begin{itemize}[leftmargin=*]
    \item credenziali (cookie, certificati);
    \item header \texttt{Content-Type} per la login e \texttt{Authorization} per il JWT;
    \item tutti i verbi HTTP rilevanti per CRUD.
\end{itemize}

\subsection{Front-end: routing principale}

Il componente \texttt{app.tsx} definisce le rotte dell'applicazione:

\begin{lstlisting}[language=tsx]
import { BooksPage } from '../features/books/books.page';
import { Routes, Route, Navigate } from 'react-router-dom';
import { LoginPage } from '../features/auth/login.page';

export function App() {
  return (
    <Routes>
      <Route path="/" element={<Navigate to="/books" replace />} />
      <Route path="/login" element={<LoginPage />} />
      <Route path="/books" element={<BooksPage />} />
    </Routes>
  );
}
export default App;
\end{lstlisting}

\subsection{API client frontend}

\subsubsection{\texttt{auth.api.ts}}
Login: invio email/password, ricezione JWT, salvataggio in \texttt{localStorage}:

\begin{lstlisting}[language=tsx]
const API_URL = 'http://localhost:3333/api';

export async function login(email: string, password: string) {
  const response = await fetch(`${API_URL}/auth/login`, {
    method: 'POST',
    headers: { 'Content-Type': 'application/json' },
    body: JSON.stringify({ email, password }),
  });

  if (!response.ok) {
    throw new Error('Credenziali non valide');
  }

  const data = await response.json();
  localStorage.setItem('access_token', data.access_token);
  return data;
}
\end{lstlisting}

\subsubsection{\texttt{books.api.ts}}
Chiamata GET autenticata: header \texttt{Authorization: Bearer <token>}:

\begin{lstlisting}[language=tsx]
import { BookListItem } from './types';
const API_URL = 'http://localhost:3333/api';

export async function fetchBooks(): Promise<BookListItem[]> {
  const token = localStorage.getItem('access_token');

  const response = await fetch(`${API_URL}/books`, {
    headers: { Authorization: `Bearer ${token}` },
  });

  if (!response.ok) {
    throw new Error('Errore durante il caricamento dei libri');
  }

  return response.json();
}
\end{lstlisting}

\subsection{Pagina di login}

\subsubsection{Stato e gestione del submit}
Usa \texttt{useState} per email, password, errore e flag di loading; \texttt{useNavigate} per il redirect:

\begin{lstlisting}[language=tsx]
export function LoginPage() {
  const [email, setEmail] = useState('');
  const [password, setPassword] = useState('');
  const [error, setError] = useState<string | null>(null);
  const [loading, setLoading] = useState(false);
  const navigate = useNavigate();

  async function handleSubmit(e: React.FormEvent) {
    e.preventDefault();
    setError(null);
    setLoading(true);

    try {
      await login(email, password);
      navigate('/books');
    } catch (err: any) {
      setError(err.message);
    } finally {
      setLoading(false);
    }
  }
}
\end{lstlisting}

\subsubsection{Rendering con CSS Module}
Il foglio di stile (\texttt{books.module.css}) \`e condiviso tra le pagine. JSX semplificato:

\begin{lstlisting}[language=tsx]
import book_styles from '../css/books.module.css';

return (
  <main className={book_styles.page}>
    <div className={book_styles.card}>
      <h1 className={book_styles.title}>Library Login</h1>
      <form onSubmit={handleSubmit} className={book_styles.form}>
        <div className={book_styles.field}>
          <label>Email</label>
          <input
            value={email}
            onChange={(e) => setEmail(e.target.value)}
            required
          />
        </div>
        <div className={book_styles.field}>
          <label>Password</label>
          <input
            type="password"
            value={password}
            onChange={(e) => setPassword(e.target.value)}
            required
          />
        </div>
        <button type="submit" disabled={loading}>
          {loading ? 'Accesso in corso...' : 'Login'}
        </button>
      </form>
      {error && <p className={book_styles.error}>{error}</p>}
    </div>
  </main>
);
\end{lstlisting}

\subsection{Pagina del catalogo}

Carica i libri al mount con \texttt{useEffect} e mostra stati di \emph{loading}, \emph{error} e \emph{empty}:

\begin{lstlisting}[language=tsx]
export function BooksPage() {
  const [books, setBooks] = useState<BookListItem[]>([]);
  const [loading, setLoading] = useState(true);
  const [error, setError] = useState<string | null>(null);
  const navigate = useNavigate();

  useEffect(() => {
    fetchBooks()
      .then(setBooks)
      .catch((err) => setError(err.message))
      .finally(() => setLoading(false));
  }, []);

  if (loading) { return <p>Caricamento libri...</p>; }
}
\end{lstlisting}

\subsubsection{Pacchetto \texttt{clsx}}
Permette di applicare pi\`u classi CSS Module insieme:

\begin{lstlisting}[language=tsx]
import clsx from 'clsx';
<section className={clsx(book_styles.card, book_styles.cardLarge)}>
\end{lstlisting}

\subsubsection{Pulsante di Logout nella header}
\begin{lstlisting}[language=tsx]
<button
  className={book_styles.logoutButton}
  onClick={() => navigate('/logout')}
>
  Logout
</button>
\end{lstlisting}

\subsubsection{Rendering della tabella}
Tabella con map sugli elementi e gestione di campi opzionali:

\begin{lstlisting}[language=tsx]
{books.length === 0 ? (
  <p>Nessun libro disponibile.</p>
) : (
  <table className={book_styles.table}>
    <thead>
      <tr>
        <th>Titolo</th><th>Anno</th>
        <th>Categoria</th><th>Autori</th>
      </tr>
    </thead>
    <tbody>
      {books.map((book) => (
        <tr key={book.id}>
          <td>{book.title}</td>
          <td>{book.publishedYear}</td>
          <td>
            <span>{book.category?.name ?? 'N/D'}</span>
          </td>
          <td>
            {book.authors?.length
              ? book.authors
                  .map(a => `${a.firstName} ${a.lastName}`)
                  .join(', ')
              : 'N/D'}
          </td>
        </tr>
      ))}
    </tbody>
  </table>
)}
\end{lstlisting}

\subsection{Componente di Logout}

Gestito interamente lato React, senza chiamate al back-end (il token JWT viene semplicemente rimosso dal \texttt{localStorage}):

\begin{lstlisting}[language=tsx]
import { useEffect } from 'react';
import { useNavigate } from 'react-router-dom';

export function LogoutPage() {
  const navigate = useNavigate();

  useEffect(() => {
    localStorage.removeItem('access_token');
    navigate('/login', { replace: true });
  }, [navigate]);

  return <p>Logout in corso...</p>;
}
\end{lstlisting}

\subsection{Proteggere una rotta}

Un componente aggiuntivo \texttt{ProtectedRoute} fa da middleware: se manca il token JWT, redirige a \texttt{/login}; altrimenti renderizza i \texttt{children}:

\begin{lstlisting}[language=tsx]
import { ReactNode } from 'react';
import { Navigate } from 'react-router-dom';

export function ProtectedRoute({ children }: { children: ReactNode }) {
  const token = localStorage.getItem('access_token');

  if (!token) {
    return <Navigate to="/login" replace />;
  }

  return children;
}
\end{lstlisting}

Si applica nell'\texttt{App} avvolgendo le rotte protette:

\begin{lstlisting}[language=tsx]
<Routes>
  <Route path="/" element={<Navigate to="/books" replace />} />
  <Route path="/login" element={<LoginPage />} />
  <Route
    path="/books"
    element={
      <ProtectedRoute>
        <BooksPage />
      </ProtectedRoute>
    }
  />
  <Route path="/logout" element={<LogoutPage />} />
</Routes>
\end{lstlisting}

% ====================================================================
\section{L8c -- Approfondimenti React (CRUD completo + framework CSS)}
% ====================================================================

\subsection{Obiettivo della lezione}

Estensione dell'esempio della Libreria online con aspetti tipici di un'applicazione web reale:
\begin{itemize}[leftmargin=*]
    \item pagina di creazione di un nuovo libro;
    \item pulsanti di modifica e cancellazione per ogni libro nel catalogo;
    \item pagina di modifica di un libro esistente;
    \item navbar condivisa.
\end{itemize}

In parallelo si introducono i framework CSS \textbf{Tailwind} e \textbf{Bootstrap}.

\subsection{Estensioni della libreria online}

\subsubsection{Predisposizione della pagina ``Nuovo libro''}
Form con campi: Titolo, Anno pubblicazione, ID categoria (select), Autori (campo di ricerca con autocomplete + chip per gli autori selezionati). Bottone ``Crea libro''.

\subsubsection{Predisposizione di Modifica/Elimina nel catalogo}
Per ogni riga della tabella, due bottoni: \emph{Modifica} (naviga a \texttt{/books/:id/edit}) ed \emph{Elimina} (\texttt{window.confirm} + DELETE + update ottimistico della lista).

\subsubsection{Predisposizione della pagina ``Modifica libro''}
Stesso form della creazione, precompilato con \texttt{fetchBookById(id)} dentro un \texttt{useEffect}; submit con PATCH. Mostra validazioni inline tipo ``Seleziona almeno un autore''.

\subsubsection{Predisposizione delle navbar}
Componente con i link alle pagine principali (Catalogo, Nuovo libro, Autori, Categorie) e dropdown utente con logout.

\subsection{Framework CSS -- Tailwind}

\textbf{Tailwind} \`e un framework CSS \emph{utility-first}: le utility descrivono direttamente lo stile, senza dover scrivere classi semantiche personalizzate.

\begin{lstlisting}[language=tsx]
/* CSS classico */
.card {
  background: white;
  padding: 2rem;
  border-radius: 12px;
}

/* Tailwind equivalente */
<div className="bg-white p-8 rounded-xl"></div>
\end{lstlisting}

Vantaggi:
\begin{itemize}[leftmargin=*]
    \item maggiore velocit\`a di sviluppo;
    \item niente naming CSS, niente conflitti globali;
    \item design consistente;
    \item ideale per componenti React.
\end{itemize}

\subsubsection{Installazione Tailwind in Nx + React + Vite}

Tre passaggi:

\textbf{1. Installazione pacchetti.}
\begin{lstlisting}
npm install -D tailwindcss@3 postcss autoprefixer
\end{lstlisting}

\textbf{2. \texttt{apps/ui/postcss.config.cjs}.}
\begin{lstlisting}[language=tsx]
module.exports = {
  plugins: {
    tailwindcss: {},
    autoprefixer: {},
  },
};
\end{lstlisting}

\textbf{3. \texttt{apps/ui/tailwind.config.cjs}.}
\begin{lstlisting}[language=tsx]
module.exports = {
  content: ['./src/**/*.{html,js,ts,jsx,tsx}'],
  theme: { extend: {} },
  plugins: [],
};
\end{lstlisting}

\textbf{4. Direttive Tailwind in \texttt{styles.css}.}
\begin{lstlisting}
@tailwind base;
@tailwind components;
@tailwind utilities;
\end{lstlisting}

Il file \texttt{styles.css} \`e gi\`a importato in \texttt{main.tsx}.

\subsection{Framework CSS -- Bootstrap}

\textbf{Bootstrap} \`e un framework CSS che fornisce \emph{componenti gi\`a pronti} (button, card, navbar, modal, ecc.) oltre a utility classes. Pi\`u maturo e immediato di Tailwind ma meno flessibile/moderno.

Installazione minima:
\begin{lstlisting}
npm install bootstrap
\end{lstlisting}

E import nel \texttt{main.tsx}:
\begin{lstlisting}[language=tsx]
import 'bootstrap/dist/css/bootstrap.min.css';
\end{lstlisting}

\subsection{Confronto pratico: home page realizzata nei tre stili}

\subsubsection{Con CSS Modules}
La pagina importa un \texttt{styles} dall'omonimo file \texttt{.module.css} e applica le classi locali:

\begin{lstlisting}[language=tsx]
import styles from '../css/home.module.css';

export function HomeModulePage() {
  return (
    <main className={styles.page}>
      <section className={styles.card}>
        <h1 className={styles.title}>Library App</h1>
        <p className={styles.description}>
          Benvenuto nel catalogo online della biblioteca.
        </p>
        <a className={styles.button} href="/books">
          Vai al catalogo
        </a>
      </section>
    </main>
  );
}
\end{lstlisting}

\subsubsection{Con Tailwind}
Niente file CSS dedicato: tutte le utility nelle \texttt{className}.

\begin{lstlisting}[language=tsx]
import { Link } from 'react-router-dom';

export function HomeTailwindPage() {
  return (
    <main className="flex min-h-screen items-center justify-center bg-slate-100 px-4">
      <section className="w-full max-w-md rounded-2xl bg-white p-10 text-center shadow-xl">
        <h1 className="text-3xl font-bold text-slate-900">Library App</h1>
        <p className="mt-4 text-slate-500">
          Benvenuto nel catalogo online della biblioteca.
        </p>
        <Link
          to="/books"
          className="mt-8 inline-block rounded-lg bg-blue-600 px-5 py-3 font-semibold text-white hover:bg-blue-700"
        >
          Vai al catalogo
        </Link>
      </section>
    </main>
  );
}
\end{lstlisting}

\subsubsection{Con Bootstrap}
Si usano le classi standard del framework (\texttt{card}, \texttt{btn btn-primary}, \texttt{d-flex}, ecc.):

\begin{lstlisting}[language=tsx]
import { Link } from 'react-router-dom';

export function HomeBootstrapPage() {
  return (
    <main className="min-vh-100 bg-light d-flex align-items-center justify-content-center px-3">
      <section
        className="card shadow p-5 text-center"
        style={{ maxWidth: 420 }}
      >
        <h1 className="fw-bold">Library App</h1>
        <p className="text-muted mt-3">
          Benvenuto nel catalogo online della biblioteca.
        </p>
        <Link to="/books" className="btn btn-primary mt-4">
          Vai al catalogo
        </Link>
      </section>
    </main>
  );
}
\end{lstlisting}

\subsection{Link utili}
\begin{itemize}[leftmargin=*]
    \item Documentazione Tailwind: \url{https://tailwindcss.com/}
    \item Icone per React: \url{https://react-icons.github.io/react-icons/}
\end{itemize}

% ====================================================================
\section*{Conclusioni}
\addcontentsline{toc}{section}{Conclusioni}

Le quattro lezioni costituiscono il percorso completo dal modello dei dati al frontend funzionante:

\begin{itemize}[leftmargin=*]
    \item \textbf{L7} fornisce il metodo (\textbf{WebML}) per derivare modello dati, gruppi utente e casi d'uso dai requisiti di business; questo \`e il punto di partenza che guida sia il backend sia il frontend.
    \item \textbf{L8} introduce \textbf{React}: SPA, componenti funzionali, JSX, state, eventi, rendering condizionale, liste, stili e Virtual DOM.
    \item \textbf{L8b} applica concretamente React costruendo un'app di catalogo libri completa di \textbf{autenticazione JWT}, \textbf{routing protetto}, comunicazione \textbf{CORS} con NestJS e logout lato client.
    \item \textbf{L8c} estende l'app con \textbf{CRUD completo} (creazione, modifica, eliminazione) e mostra a confronto tre approcci di stile: \textbf{CSS Modules}, \textbf{Tailwind} e \textbf{Bootstrap}.
\end{itemize}

Lo schema architetturale risultante \`e quello tipico delle SPA moderne: API REST autenticata via JWT (backend NestJS) + SPA React Router (frontend), con DTO condivisi tra le due tier per garantire coerenza dei tipi.

\end{document}
